Vysoká zaměstnanost a~dostupnost kvalifikované pracovní síly za~mzdy výrazně pod úrovní sousedních ekonomik vytvářejí z~\aclgen{ČR} atraktivní destinaci pro podniky se zahraniční kontrolou. Mapa zachycuje podíl zaměstnanců pracujících v~takovýchto podnicích na~celkové zaměstnanosti; \ac{ČR} se řadí k~zemím s~nadprůměrnou hodnotou tohoto ukazatele v~rámci \acs{geo-EU27}.

Vysoký podíl zahraniční kontroly v~podnikatelské sféře je z~hlediska sociálního dialogu ambivalentní. Na~jedné straně přitahuje kapitál, technologie a~zaměstnanost; na~druhé straně nadnárodní podniky zpravidla vyvíjejí silnější odpor vůči kolektivnímu vyjednávání na~úrovni odvětví a~preferují podnikové nebo individuální mzdové nastavení, které snižuje koordinační roli odborů. Srovnání s~\aclins{AT} nebo \aclins{DK}, kde je podíl zahraniční kontroly srovnatelný nebo vyšší, ale pokrytí \aclins{KS} zároveň vysoké, ukazuje, že přítomnost nadnárodního kapitálu sama o~sobě kolektivní vyjednávání neblokuje; rozhodující je institucionální rámec, v~němž tyto podniky operují.

\inputpgffigure{eu_mne_mapa}
